% SGA 3，第 VIII 篇，英文重建底本说明。
% 范围：第 5 节全文，止于第 6 节标题之前。
% 权威源：Polo--Gille Exp8-8nov09.pdf, 单篇第 14--17 页
% (合阅本第 630--633 页).
% 英文底本说明：Jacob Reinhold 的 Markdown 传承版本仅用于措辞对照。
% 正文、公式、编号、引用与注释以 Polo--Gille PDF 为准。

\setcounter{footnote}{16}

% 权威源单篇第 14 页; 原版页码标记 15.
\SGARefTarget{sga3:viii:section:5}\section*{5. 仿射概形对自由作用的可对角化群的商}
\addcontentsline{toc}{section}{5. 仿射概形对自由作用的可对角化群的商}

\noindent\footnote{编者注：增加了下句。}
用 \(\mathcal M\) 表示忠实平坦拟紧态射的集合，并约定挠子均指关于 \((\mathrm{fpqc})\) 拓扑的挠子。

\medskip
\noindent\SGARefTarget{sga3:viii:theorem:5:1}\textbf{定理 5.1.}
\begin{itshape}
设 \(S\) 为预概形，\(M\) 为通常的交换群，\(G=D_S(M)\) 为它定义的 \(S\) 上的可对角化群，\(P\) 为在 \(S\) 上仿射的 \(S\)-预概形，且 \(G\) 在其上自由地右作用。

则 \(G\) 在 \(P\) 上定义的等价关系为 \(\mathcal M\)-有效的（参见 IV，\SGARefLink{sga3:iv:subsection:3:4}{3.4}）；即商 \(X=P/G\) 存在，且 \(P\) 为 \(X\) 上 \(G_X=D_X(M)\) 下的挠子。此外，\(P/G\) 在 \(S\) 上仿射。更确切地说，若 \(P\) 由 \(M\)-分次代数 \(\mathcal A\) 定义，则 \(P/G\) 同构于 \(\operatorname{Spec}(\mathcal A_0)\)，其中 \(\mathcal A_0=\mathcal A^G\) 为 \(\mathcal A\) 的零次分量。
\end{itshape}

\smallskip
\noindent\textit{证明。}
置 \(X=\operatorname{Spec}(\mathcal A_0)\)。包含映射 \(\mathcal A_0\to\mathcal A\) 给出在 \(G\) 的作用下不变的自然态射 \(P\to X\)。于是 \(P\) 成为在 \(X\) 上仿射且带有算子群 \(G_X=D_X(M)\) 的 \(X\)-预概形；\(G\) 在 \(P/S\) 上自由作用的假设蕴含 \(G_X\) 在 \(P/X\) 上自由作用。利用 \(\mathcal B_0=\mathcal O_X\) 这一事实（其中 \(\mathcal B\) 为定义 \(P/X\) 的 \(X\) 上的 \(M\)-分次代数），只须证明 \(P\) 为 \(G_X\) 下的主齐性丛。

因此，可以假设 \(X=S\) 且 \(S\) 仿射。此时 \(P\) 仿射，由 \(M\)-分次环 \(A\) 定义；用 \(A_m\) 表示其 \(m\) 次齐次部分，于是 \(S=\operatorname{Spec}(A_0)\)。根据 \SGARefLink{sga3:viii:proposition:4:6}{4.6}，只需验证
\[
  A_m\cdot A_{-m}=A_0
  \qquad\text{对每个 \(m\in M\)。}
  \tag{x}\label{sga3:viii:theorem:5:1:eq:x:01}
\]
直接计算也表明，\SGARefLink{sga3:viii:theorem:5:1:eq:x:01}{(x)} 等价于说
\[
  P\times_SG\longrightarrow P\times_SP
\]
为闭浸入，而不仅仅是单态射（后者来自 \(G\) 自由作用的假设）；等价地，仿射环的同态

% 权威源单篇第 15 页; 原版页码标记 16.
\[
  \theta:A\otimes_{A_0}A\longrightarrow A(M)
\]
为满射。\footnote{编者注：事实上，对每个 \(b\in A\) 和 \(a_m\in A_m\)，有 \(\theta(b\otimes a_m)=ba_m\otimes e_m\)；因此 \(\theta\) 的满射性等价于 \SGARefLink{sga3:viii:theorem:5:1:eq:x:01}{(x)}。}
若假设 \(G\) 在 \(P\) 上定义的等价关系是闭的，这就证明了 \SGARefLink{sga3:viii:theorem:5:1}{5.1}。不过，我们将证明此假设已可由 \(G\) 自由作用推出。这也隐含于\SGARefLink{sga3:viii:theorem:5:1}{定理 5.1} 中：概形 \(G\times_SP\) 必同构于 \(P\times_XP\)，而后者在 \(P\times_SP\) 中是闭的，因为 \(X\) 在 \(S\) 上仿射，从而分离。

令 \(R=P\times_SG\)。\(G\) 自由作用，即 \(R\to P\times_SP\) 为单态射的假设，意味着对角态射
\[
  R\longrightarrow
  R\times_{(P\times_SP)}R=R'
\]
为同构。有 \(R=\operatorname{Spec}(A(M))\) 以及
\[
  R'=\operatorname{Spec}\bigl(A(M\times M)/K\bigr),
\]
\footnote{编者注：将原文记为 \(J\) 的理想记为 \(K\)，以区别于 \SGARefLink{sga3:viii:theorem:5:1:eq:xxx:01}{(xxx)} 中 \(A_0\) 的理想 \(J_p\)。下文还明确指出，关系 \SGARefLink{sga3:viii:theorem:5:1:eq:xx:01}{(xx)} 等价于等式 \(K=K'\)；见后文。}
其中 \(K\) 为由元素
\[
  x_m(e_{m,0}-e_{0,m}),
  \qquad m\in M,\quad x_m\in A_m.
\]
生成的理想。令
\[
  \phi:A(M\times M)\longrightarrow A(M)
\]
为由
\[
  x e_{m,n}\longmapsto x e_{m+n}
  \qquad (m,n\in M,\ x\in A),
\]
定义的满环同态，其中各 \(e_m\)、各 \(e_{m,n}=e_m\otimes e_n\) 分别为 \(A(M)\)、\(A(M\times M)\) 的典范基的元素。于是对角态射 \(R\to R'\) 对应于模去 \(K\) 后得到的同态
\[
  \bar\phi:A(M\times M)/K\longrightarrow A(M)
\]
\(\phi\) 的核为由
\[
  d_m=e_{m,0}-e_{0,m}.
\]
生成的理想 \(K'\)。
\emph{〔原文印作 \(\ker(\bar\phi)=K'\)；模去 \(K\) 后，相应的核为 \(K'/K\)。〕}
有 \(K\subseteq K'\)，而 \(G\) 在 \(P\) 上自由作用，即 \(\bar\phi\) 为同构的假设，等价于 \(K'=K\)，可表为关系
\[
  d_m\in K=\sum_p A(M\times M)\,A_p\,d_p
  \qquad\text{对每个 \(m\in M\)。}
  \tag{xx}\label{sga3:viii:theorem:5:1:eq:xx:01}
\]

利用 \(A(M\times M)\) 的自然三重分次，以及 \(d_m\) 的第一个次数为零这一事实，这就是说 \(d_m\) 可以写成形如
\[
  f e_{r,s}(e_{p,0}-e_{0,p}),
  \qquad f\in A_{-p}\cdot A_p.
\]
的元素之和。由于 \(d_m\) 的总次数为 \(m\)，只需采用满足
\[
  r+s+p=m.
\]
的项。

% 权威源单篇第 16 页; 原版页码标记 17.
因此，对每个 \(m\in M\)，必有表达式
\[
  \left\{
  \begin{aligned}
    d_m=e_{m,0}-e_{0,m}
      &=\sum_{r,s}\lambda_{r,s}
        (e_{m-s,s}-e_{r,m-r}),\\
    \lambda_{r,s}&\in J_p=A_p\cdot A_{-p}\subseteq A_0,
      \qquad p=m-(r+s).
  \end{aligned}
  \right.
  \tag{xxx}\label{sga3:viii:theorem:5:1:eq:xxx:01}
\]
必须由此推出 \SGARefLink{sga3:viii:theorem:5:1:eq:x:01}{(x)}，即关系
\[
  J_n=A_0
  \qquad\text{对每个 \(n\in M\)。}
  \tag{xxxx}\label{sga3:viii:theorem:5:1:eq:xxxx:01}
\]
为此，只须模 \(A_0\) 的每个极大理想证明\SGARefLink{sga3:viii:theorem:5:1:eq:xxxx:01}{同一关系}。由于所有假设在这样的约化下保持不变，可以假设 \(A_0\) 为域。

\medskip
\noindent\SGARefTarget{sga3:viii:lemma:5:2}\textbf{引理 5.2.}
\textit{在上述条件下，设 \(A_0\) 为域。若 \(M\ne0\)，则存在 \(p\in M-\{0\}\)，使 \(J_p=A_0\)。}

事实上，否则 \SGARefLink{sga3:viii:theorem:5:1:eq:xxx:01}{(xxx)} 右边的和对每个 \(m\in M\) 都为零，这是荒谬的。

% 原版页码标记 18.
\medskip
\noindent\SGARefTarget{sga3:viii:corollary:5:3}\textbf{推论 5.3.}
\textit{在上述条件下，但不假设 \(A_0\) 为域，则存在有限多个元素 \(p_i\in M-\{0\}\)，使}
\[
  \sum_iJ_{p_i}=A_0.
\]

事实上，将 \SGARefLink{sga3:viii:lemma:5:2}{5.2} 用于由 \(A/A_0\) 模 \(A_0\) 的极大理想约化所得的情形即可。\footnote{编者注：由 \SGARefLink{sga3:viii:lemma:5:2}{5.2} 得 \(\sum_{p\ne0}J_p=A_0\)；因此 \(1\) 可写为有限和 \(\sum_i x_i\)，其中 \(x_i\in J_{p_i}\)。}

\medskip
\noindent\SGARefTarget{sga3:viii:corollary:5:4}\textbf{推论 5.4.}
\textit{再设 \(A_0\) 为域。对 \(M\) 的每个满足 \(N\ne M\) 的子群 \(N\)，存在 \(p\in M-N\)，使 \(J_p=A_0\)。}

事实上，令 \(M'=M/N\)，考虑 \(M'\)-分次环 \(A'\)，其底环为 \(A\)，分次由
\[
  A'_{m'}=\coprod_{m\in h^{-1}(m')}A_m,
\]
给出，其中 \(h:M\to M'=M/N\) 为典范同态。在几何上，这一构造相当于考虑 \(G\) 的子群 \(G'=D_S(M')\)，以及由 \(G\) 的作用诱导的 \(P\) 上带有算子群 \(G'\) 的概形结构。于是显然 \(G'\) 在 \(P\) 上自由作用\emph{〔原文印作 \(P'\)〕}；即对 \((M',A')\) 满足 \SGARefLink{sga3:viii:corollary:5:3}{5.3} 的假设。因此
\[
  1=\sum_i f_i g_i,
  \qquad
  f_i\in A'_{m_i'},\quad
  g_i\in A'_{-m_i'},\quad
  m_i'\in M'-\{0\}.
\]
取右边沿 \(A_0\) 的分量，并利用 \(A_0\) 为域这一事实，立即得到 \SGARefLink{sga3:viii:corollary:5:4}{5.4}。

现在注意到
\[
  J_pJ_q\subseteq J_{p+q}
  \qquad\text{及}\qquad
  J_p=J_{-p}.
\]

% 权威源单篇第 17 页; 原版页码标记 19.
因此，若 \(N\) 表示满足 \(J_m=A_0\) 的 \(m\in M\) 的集合\emph{〔原文印作 \(J_p=A_0\)〕}，则 \(N\) 为 \(M\) 的子群。由 \SGARefLink{sga3:viii:corollary:5:4}{5.4}，它等于 \(M\)。这就完成了\SGARefLink{sga3:viii:theorem:5:1}{定理 5.1} 的证明。

如证明过程中所指出的，\SGARefLink{sga3:viii:theorem:5:1}{定理 5.1} 蕴含：

\medskip
\noindent\SGARefTarget{sga3:viii:corollary:5:5}\textbf{推论 5.5.}
\textit{在 \SGARefLink{sga3:viii:theorem:5:1}{5.1} 的假设下，图态射}
\[
  P\times_SG\longrightarrow P\times_SP
\]
\textit{为闭浸入。}

立即可得：

\medskip
\noindent\SGARefTarget{sga3:viii:corollary:5:6}\textbf{推论 5.6.}
\textit{设 \(\sigma\) 为 \(P\) 在 \(S\) 上的截面。则由 \(\sigma\) 定义的从 \(G\) 到 \(P\) 的态射 \(g\mapsto\sigma\cdot g\) 为闭浸入。}

\medskip
\noindent\SGARefTarget{sga3:viii:corollary:5:7}\textbf{推论 5.7.}
\textit{设 \(G\) 和 \(H\) 为两个 \(S\)-群，其中 \(G\) 可对角化，\(H\) 在 \(S\) 上仿射；设 \(u:G\to H\) 为 \(S\)-群的同态，且为单态射。则 \(u\) 为闭浸入，\(H/G=X\) 存在，\(H\) 为 \(X\) 上 \(G_X\) 下的主齐性丛，且 \(X\) 在 \(S\) 上仿射。}

\medskip
\noindent\SGARefTarget{sga3:viii:corollary:5:8}\textbf{推论 5.8.}
\textit{在 \SGARefLink{sga3:viii:theorem:5:1}{5.1} 的假设下，若 \(P\) 在 \(S\) 上有限型（分别为有限呈示），则 \(X=P/G\) 也如此。}

事实上，假设蕴含 \(G_X\) 的纤维有限型。因此，由 \SGARefLink{sga3:viii:proposition:2:1:item:b:01}{2.1(b)}，\(G_X\) 在 \(X\) 上有限呈示；由于 \(P\) 为 \(G_X\) 下的挠子，它在 \(X\) 上也有限呈示。\footnote{编者注：参见 EGA II，2.7.1(vi)。}
它还在 \(X\) 上忠实平坦，所以结论可由第 V 篇，命题 \SGARefLink{sga3:v:proposition:9:1}{9.1} 得出。
